\documentclass[11pt,oneside]{book}

\usepackage[letterpaper,margin=1.05in]{geometry}
\usepackage{amsmath,amssymb,amsthm,mathtools}
\usepackage{booktabs,longtable,array}
\usepackage{enumitem}
\usepackage{xcolor}
\usepackage{hyperref}

\hypersetup{
  colorlinks=true,
  linkcolor=blue!45!black,
  urlcolor=blue!45!black,
  citecolor=blue!45!black
}

\setlength{\parindent}{0pt}
\setlength{\parskip}{6pt}
\setlist[itemize]{leftmargin=1.5em}
\setlist[enumerate]{leftmargin=1.8em}
\sloppy

\theoremstyle{definition}
\newtheorem{definition}{Definition}[chapter]
\newtheorem{theorem}{Theorem}[chapter]
\newtheorem{proposition}{Proposition}[chapter]
\newtheorem{remark}{Remark}[chapter]

\newenvironment{proofsteps}
{\begin{enumerate}[label=\textbf{Line \arabic*.}, itemsep=0.35em]}
{\end{enumerate}}

\newcommand{\E}{\mathbb{E}}
\newcommand{\Q}{\mathbb{Q}}
\newcommand{\R}{\mathbb{R}}
\newcommand{\Var}{\operatorname{Var}}
\newcommand{\Cov}{\operatorname{Cov}}
\newcommand{\one}{\mathbf{1}}
\newcommand{\dd}{\,\mathrm{d}}
\newcommand{\clip}{\operatorname{clip}}
\newcommand{\se}{\operatorname{SE}}
\newcommand{\argmin}{\operatorname*{arg\,min}}

\title{\textbf{Monte Carlo Proxy Pricing}\\
\large A textbook-style methodology for European, American, Asian, barrier, cliquet, basket Asian, SLV cliquet, basket cliquet, autocallable, and random-payoff proxies}
\author{Proxy Pricing Research Project}
\date{\today}

\begin{document}
\frontmatter
\maketitle
\tableofcontents

\chapter*{Reader's Guide}
This document is written as a second-year undergraduate text. It assumes calculus, linear algebra, basic probability, and a willingness to read a few stochastic-process formulas. The goal is not to hide the implementation behind jargon. The goal is to show exactly how a Monte Carlo price label becomes a fast proxy.

The main chapters explain the workflow. The appendices give line-by-line proofs for the formulas used in the code. Every nontrivial implementation trick is classified as one of:
\begin{itemize}
  \item an exact pricing identity,
  \item an unbiased Monte Carlo variance-reduction identity,
  \item a numerical approximation whose error must be validated.
\end{itemize}

\mainmatter

\chapter{The Pricing Problem}

\section{Risk-neutral value}
Let \(H\) be the payoff of a derivative at maturity or final settlement time \(T\). Under the risk-neutral measure \(\Q\), the time-\(t\) value is
\[
  V_t = B_t\,\E_\Q\!\left[\frac{H}{B_T}\,\middle|\,\mathcal F_t\right],
\]
where \(B_t=\exp(\int_0^t r_s\dd s)\) is the money-market account. If the short rate is constant,
\[
  V_t = e^{-r(T-t)}\E_\Q[H\mid\mathcal F_t].
\]

In the code, the expensive object is the conditional expectation. The proxy is a fast approximation of the map
\[
  x \longmapsto V(t,x),
\]
where \(x\) is the state available at the valuation date.

\section{What makes a good proxy state}
A state is good when it is sufficient, low-dimensional, and scaled sensibly.

\begin{definition}[Sufficient state]
A vector \(X_t\) is sufficient for pricing if, after conditioning on \(X_t\), the distribution of all future cash flows no longer depends on earlier path history.
\end{definition}

Examples:
\begin{itemize}
  \item European option: \(X_t=S_t\).
  \item Arithmetic Asian option: \(X_t=(S_t,A_t)\), where \(A_t\) is the running sum of past fixings.
  \item American put: \(X_t=S_t\), but the time index matters because of remaining exercise dates.
  \item Barrier option: \(X_t=(S_t,I_t)\), where \(I_t\) says whether the barrier has already been hit.
  \item Single-name SLV cliquet: \(X_t=(a_t,S_t,v_t)\), accrued coupon, spot, and variance.
  \item Basket SLV cliquet: \(X_t=(a_t,S_t^{(1:3)},v_t^{(1:3)})\), plus contract variant.
\end{itemize}

\section{Error decomposition}
The observed proxy error is not one number with one cause. A useful decomposition is
\[
  \widehat V_{\mathrm{proxy}} - V_{\mathrm{true}}
  =
  \underbrace{\widehat V_{\mathrm{proxy}} - V_{\mathrm{MC-label}}}_{\text{fitting error}}
  +
  \underbrace{V_{\mathrm{MC-label}}-V_{\mathrm{model}}}_{\text{training MC noise}}
  +
  \underbrace{V_{\mathrm{model}}-V_{\mathrm{true}}}_{\text{model/discretization error}}.
\]
The experiments in this repository mostly study the first two terms under fixed model assumptions.

\chapter{Monte Carlo Label Construction}

\section{Plain estimator}
For a state \(x\), define the discounted path value
\[
  Y(x,Z)=e^{-r\tau}H(x,Z),
\]
where \(Z\) represents all future random shocks. The target label is
\[
  V(x)=\E[Y(x,Z)].
\]
With \(N\) simulated paths,
\[
  \widehat V_N(x)=\frac1N\sum_{n=1}^N Y(x,Z_n).
\]
The estimator is unbiased when the paths have the correct law. Its standard error is estimated by
\[
  \widehat{\se}(\widehat V_N)
  =
  \sqrt{\frac{1}{N(N-1)}\sum_{n=1}^N\bigl(Y_n-\overline Y\bigr)^2}.
\]

\section{Sobol and antithetics}
Sobol points reduce integration error by filling the unit cube more evenly than independent random points. The implementation rounds the number of paths to a power of two because Sobol nets are naturally balanced in blocks of \(2^m\).

Antithetic pairing uses the fact that, for a standard normal vector \(Z\),
\[
  Z \stackrel{d}= -Z.
\]
The paired payoff
\[
  \frac{f(Z)+f(-Z)}{2}
\]
has the same expectation as \(f(Z)\). It often has lower variance for monotone or nearly monotone payoffs.

\section{Likelihood-ratio importance sampling}
State enrichment and importance sampling are different.

State enrichment changes the distribution of training states \(x_i\). It does not change the path law for any individual conditional expectation.

Importance sampling changes the path density from \(p\) to \(q\) and multiplies by a likelihood ratio:
\[
  \E_p[h(Z)]
  =
  \E_q\!\left[h(Z)\frac{p(Z)}{q(Z)}\right].
\]
For a normal mean shift, \(q(z)=\phi(z-\theta)\), so
\[
  \frac{p(z)}{q(z)}
  =
  \exp\!\left(-\theta z+\frac12\theta^2\right).
\]
This identity is exact; its usefulness depends on variance.

\section{Control variates}
If \(C\) has known mean \(m_C\), then
\[
  Y_\beta = Y-\beta(C-m_C)
\]
has the same expectation as \(Y\). The variance-minimizing coefficient is
\[
  \beta^\star=\frac{\Cov(Y,C)}{\Var(C)}.
\]
Asian options use a geometric Asian control when the arithmetic payoff and geometric payoff are strongly correlated.

\chapter{Proxy Fitting Principles}

\section{Targets}
Raw prices can be hard to fit because they range from nearly zero to almost linear. The code uses transformations such as:
\[
  \log(V+\varepsilon),\qquad
  \log\frac{V+\varepsilon}{B+\varepsilon},\qquad
  \operatorname{logit}\frac{V-L}{U-L}.
\]
Here \(B\) is a baseline, and \(L,U\) are known lower and upper bounds.

\section{Exact regions}
Never ask a regressor to learn something known analytically.

Examples:
\begin{itemize}
  \item At maturity, use the payoff exactly.
  \item Deep Asian call in-the-money region: the positive part disappears and the value is linear.
  \item Cliquet global floor/cap tails: the final clipped payoff is already known.
  \item American/Bermudan exercise: impose \(V=\max(\text{intrinsic},\text{continuation})\).
\end{itemize}

\section{Relative error floor}
The reported relative error is
\[
  e_{\mathrm{rel}}
  =
  \frac{|\widehat V-V|}{\max(|V|,\eta)}.
\]
This is ordinary relative error when \(|V|\ge \eta\), and it avoids meaningless exploding percentages when \(V\) is tiny.

\chapter{One-Dimensional Shape-Preserving Proxies}

\section{Why PCHIP is the default for 1D}
For a single feature, PCHIP is a strong operational default:
\begin{itemize}
  \item it interpolates labels exactly;
  \item it is local, so one noisy point does not create global oscillation;
  \item it preserves monotone shape when the data are monotone;
  \item it behaves well near exercise kinks, barrier transitions, and low-value tails.
\end{itemize}

\section{Cubic Hermite form}
On \([x_i,x_{i+1}]\), let \(h_i=x_{i+1}-x_i\), \(t=(x-x_i)/h_i\). A cubic Hermite segment is
\[
\begin{aligned}
H_i(t)
&=y_i(2t^3-3t^2+1)
  +h_i d_i(t^3-2t^2+t)\\
&\quad
  +y_{i+1}(-2t^3+3t^2)
  +h_i d_{i+1}(t^3-t^2).
\end{aligned}
\]
PCHIP chooses slopes \(d_i\) from neighboring secants in a way that prevents artificial extrema.

\section{Akima, splines, and Bernstein}
Akima interpolation is also local and robust to outliers. Natural cubic splines are smoother but can overshoot. Bernstein polynomials are shape-friendly on bounded intervals and were competitive for some barrier variants. For one-dimensional production defaults, the current rule is:
\[
  \text{PCHIP first, compare Akima/Bernstein when the validation surface says so.}
\]

\chapter{Sparse Chebyshev and Higher Dimension}

\section{Chebyshev scaling}
Chebyshev polynomials satisfy
\[
  T_n(\cos\theta)=\cos(n\theta),
\]
so \(|T_n(x)|\le 1\) for \(x\in[-1,1]\). This is why features are scaled before fitting:
\[
  z_j = 2\frac{x_j-\ell_j}{u_j-\ell_j}-1.
\]

\section{Ridge regression}
Given a design matrix \(A\), labels \(y\), and penalty matrix \(\Gamma\), ridge regression solves
\[
  \min_\beta \|A\beta-y\|_2^2+\lambda\|\Gamma\beta\|_2^2.
\]
The normal equation is
\[
  (A^\top A+\lambda\Gamma^\top\Gamma)\beta=A^\top y.
\]

\section{Why sparsity is necessary}
A full tensor degree-\(p\) basis in \(D\) variables contains \((p+1)^D\) terms. With \(D=7\) and \(p=5\), that is \(279{,}936\) terms before fitting. Basket cliquet therefore needs sparse terms, payoff-aware features, or a slower safety method.

\chapter{Instrument Recipes}

\section{European option}
State: \(S\). Coordinate:
\[
  d_1=\frac{\log(S/K)+(r-q+\tfrac12\sigma^2)\tau}{\sigma\sqrt{\tau}}.
\]
Training labels use shifted Sobol terminal GBM paths. Validation uses Black-Scholes:
\[
  C=S e^{-q\tau}\Phi(d_1)-K e^{-r\tau}\Phi(d_2),\qquad d_2=d_1-\sigma\sqrt{\tau}.
\]
Default proxy: log price fitted by PCHIP in \(d_1\).

\section{Arithmetic Asian option}
State: \(S\) and running sum \(A\). With \(m\) future fixings and total \(N\),
\[
  K_{\mathrm{adj}}=\frac{NK-A-S}{m}.
\]
The call payoff can be rewritten as
\[
  \left(\frac{A+S+\sum_{j=1}^m S G_j}{N}-K\right)_+
  =
  S\frac{m}{N}\left(\frac1m\sum_{j=1}^mG_j-\frac{K_{\mathrm{adj}}}{S}\right)_+.
\]
Thus the tested monthly call can be represented by an adjusted-moneyness coordinate. Training uses Sobol antithetics, likelihood shifts where useful, and a geometric Asian control variate.

\section{Ten-asset basket Asian}
State: ten spots plus the running basket sum. The proxy starts from a moment-matched lognormal baseline for the final arithmetic basket average, then fits a sparse Chebyshev/PCA correction and a PCHIP residual calibration. PCA summarizes cross-sectional composition without throwing away the level and running-sum information.

\section{American and Bermudan options}
At each exercise date \(k\),
\[
  V_k(S)=\max\!\left(g_k(S),\,e^{-r\Delta t}\E[V_{k+1}(S_{k+1})\mid S_k=S]\right).
\]
The continuation value is estimated by one-step Sobol transitions and fitted with PCHIP. The max with intrinsic value is imposed exactly.

\section{Barrier option}
For a zero-rebate barrier,
\[
  V_{\mathrm{in}}+V_{\mathrm{out}}=V_{\mathrm{vanilla}}.
\]
Discrete monitoring checks simulated observation dates. Continuous monitoring uses Brownian-bridge survival between observations. For a lower log barrier \(b\), endpoints \(x,y>b\), and variance \(\sigma^2\Delta t\),
\[
  P(\text{survive}\mid x,y)=1-\exp\!\left(-\frac{2(x-b)(y-b)}{\sigma^2\Delta t}\right).
\]

\section{Cliquet and SLV cliquet}
Cliquet payoffs are bounded by local and global caps/floors. If accrued value \(a\) and \(m\) remaining coupons imply the final total is certainly below the global floor or above the global cap, the value is exact.

For single-name SLV cliquet, the state is \((a,S,v)\). The proxy uses exact tails, bounded-logit targets, and local/sparse Chebyshev features. For three-asset basket cliquet, the fitted proxy is not universally enough for strict max-error control; hard order-statistic cases use a cached Sobol/LR safety proxy.

\section{Autocallable and random payoff}
Autocallable valuation is naturally date-indexed. The state is spot at an observation date; the cash-flow rule determines whether the note redeems early or continues. Random payoff tests are synthetic one-dimensional payoffs used to stress interpolation under different shapes.

\chapter{Validation and Timing}

\section{Independent benchmark principle}
Training labels and benchmark labels use separate random streams. Benchmarks use more paths, closed forms, finite differences, or tree methods when available.

\section{Current reporting}
The timing report separates:
\begin{itemize}
  \item training sample generation;
  \item proxy fitting after labels are available;
  \item option-family wall time;
  \item average wall time per scenario.
\end{itemize}
The reduced-budget smoke accuracy columns in the timing report are diagnostic only. The product-specific markdown summaries remain the full-budget accuracy source.

\appendix

\chapter{Line-by-Line Proofs: Monte Carlo}

\section{Unbiased sample mean}
\begin{proofsteps}
\item Define \(Y=e^{-r\tau}H\) and \(V=\E_\Q[Y]\).
\item Define \(\widehat V_N=N^{-1}\sum_{n=1}^N Y_n\).
\item Apply linearity:
\[
\E[\widehat V_N]
=\E\!\left[\frac1N\sum_{n=1}^N Y_n\right]
=\frac1N\sum_{n=1}^N\E[Y_n]
=\frac1N\sum_{n=1}^N V
=V.
\]
\item If paths are independent with variance \(\sigma_Y^2\),
\[
\Var(\widehat V_N)
=\Var\!\left(\frac1N\sum_{n=1}^N Y_n\right)
=\frac1{N^2}\sum_{n=1}^N\Var(Y_n)
=\frac{\sigma_Y^2}{N}.
\]
\item Replace \(\sigma_Y^2\) by the sample variance to estimate standard error.
\end{proofsteps}

\section{Likelihood-ratio shift}
\begin{proofsteps}
\item Start with \(V=\int h(z)p(z)\dd z\).
\item Choose a proposal \(q\) with \(q(z)>0\) wherever \(h(z)p(z)\neq 0\).
\item Multiply and divide by \(q\):
\[
V=\int h(z)\frac{p(z)}{q(z)}q(z)\dd z
=\E_q[h(Z)L(Z)],\qquad L(Z)=\frac{p(Z)}{q(Z)}.
\]
\item For \(p(z)=\phi(z)\), \(q(z)=\phi(z-\theta)\),
\[
L(z)
=\frac{\exp(-z^2/2)}{\exp(-(z-\theta)^2/2)}
=\exp\!\left(-\theta z+\frac{\theta^2}{2}\right).
\]
\item For independent shocks, likelihood ratios multiply:
\[
L(z_1,\ldots,z_m)
=\prod_{j=1}^m\exp\!\left(-\theta_j z_j+\frac{\theta_j^2}{2}\right).
\]
\end{proofsteps}

\section{Antithetic estimator}
\begin{proofsteps}
\item Standard normal symmetry gives \(Z\stackrel d=-Z\).
\item Define \(A=(f(Z)+f(-Z))/2\).
\item Then
\[
\E[A]=\frac{\E[f(Z)]+\E[f(-Z)]}{2}
=\frac{\E[f(Z)]+\E[f(Z)]}{2}
=\E[f(Z)].
\]
\item The variance is
\[
\Var(A)=\frac12\Var(f(Z))+\frac12\Cov(f(Z),f(-Z)).
\]
\item The pairing helps whenever the covariance term is negative.
\end{proofsteps}

\chapter{Line-by-Line Proofs: GBM and European Options}

\section{GBM solution}
\begin{proofsteps}
\item Start from \(\dd S_t/S_t=(r-q)\dd t+\sigma\dd W_t\).
\item Apply Ito's lemma to \(f(S)=\log S\):
\[
\dd\log S_t
=\frac1{S_t}\dd S_t-\frac12\frac1{S_t^2}(\dd S_t)^2.
\]
\item Since \((\dd S_t)^2=\sigma^2S_t^2\dd t\),
\[
\dd\log S_t=(r-q-\tfrac12\sigma^2)\dd t+\sigma\dd W_t.
\]
\item Integrate over \([t,T]\):
\[
\log(S_T/S_t)=(r-q-\tfrac12\sigma^2)\tau+\sigma(W_T-W_t).
\]
\item Use \(W_T-W_t=\sqrt{\tau}Z\), \(Z\sim N(0,1)\), and exponentiate.
\end{proofsteps}

\section{Black-Scholes call}
\begin{proofsteps}
\item Risk-neutral pricing:
\[
C=e^{-r\tau}\E[(S_T-K)_+].
\]
\item Split the positive part:
\[
\E[(S_T-K)_+]
=\E[S_T\mathbf 1_{\{S_T>K\}}]-K\Q(S_T>K).
\]
\item The event \(S_T>K\) is \(Z>-d_2\), so \(\Q(S_T>K)=\Phi(d_2)\).
\item Completing the square gives
\[
\E[S_T\mathbf 1_{\{S_T>K\}}]
=S e^{(r-q)\tau}\Phi(d_1).
\]
\item Substitute and discount:
\[
C=S e^{-q\tau}\Phi(d_1)-K e^{-r\tau}\Phi(d_2).
\]
\end{proofsteps}

\chapter{Line-by-Line Proofs: Asian State Reduction}

\section{Adjusted moneyness}
\begin{proofsteps}
\item At a valuation date, write future fixings as \(S G_1,\ldots,S G_m\).
\item The arithmetic average is
\[
\bar S=\frac{A+S+S\sum_{j=1}^mG_j}{N}.
\]
\item Define
\[
K_{\rm adj}=\frac{NK-A-S}{m}.
\]
\item Rewrite the call payoff:
\[
(\bar S-K)_+
=\frac1N\left(A+S+S\sum_{j=1}^mG_j-NK\right)_+
=\frac{m}{N}\left(S\frac1m\sum_{j=1}^mG_j-K_{\rm adj}\right)_+.
\]
\item Factor out \(S>0\):
\[
(\bar S-K)_+
=S\frac{m}{N}\left(\frac1m\sum_{j=1}^mG_j-\frac{K_{\rm adj}}{S}\right)_+.
\]
\item Under GBM the growth factors \(G_j\) do not depend on the current spot level \(S\).
\item Therefore \(V(S,A)=S F(K_{\rm adj}/S,t)\), which justifies a one-dimensional adjusted-moneyness proxy for the tested call setup.
\end{proofsteps}

\chapter{Line-by-Line Proofs: Barriers}

\section{In/out parity}
\begin{proofsteps}
\item Split all paths into two disjoint events: hit and no hit.
\item On hit paths, the knock-in pays the vanilla payoff and the knock-out pays zero.
\item On no-hit paths, the knock-out pays the vanilla payoff and the knock-in pays zero.
\item Therefore, path by path,
\[
H_{\rm in}+H_{\rm out}=H_{\rm vanilla}.
\]
\item Discount and take expectations:
\[
V_{\rm in}+V_{\rm out}=V_{\rm vanilla}.
\]
\end{proofsteps}

\section{Brownian-bridge lower-barrier survival}
\begin{proofsteps}
\item Condition on log endpoints \(X_0=x\), \(X_{\Delta t}=y\), both above lower barrier \(b\).
\item By the reflection principle, the conditional probability of crossing below \(b\) is
\[
\exp\!\left(-\frac{2(x-b)(y-b)}{\sigma^2\Delta t}\right).
\]
\item Survival is one minus crossing:
\[
P_{\rm surv}=1-\exp\!\left(-\frac{2(x-b)(y-b)}{\sigma^2\Delta t}\right).
\]
\item Conditional survival over many intervals is the product of interval survival probabilities.
\end{proofsteps}

\chapter{Line-by-Line Proofs: Exercise and Cliquet Bounds}

\section{American recursion}
\begin{proofsteps}
\item At final date \(M\), \(V_M(S)=g_M(S)\).
\item Suppose \(V_{k+1}\) is known.
\item Continuing from date \(k\) has value
\[
C_k(S)=e^{-r\Delta t}\E[V_{k+1}(S_{k+1})\mid S_k=S].
\]
\item Stopping immediately has value \(g_k(S)\).
\item The holder chooses the larger:
\[
V_k(S)=\max(g_k(S),C_k(S)).
\]
\item Backward induction proves optimality at every date.
\end{proofsteps}

\section{Cliquet exact tails}
\begin{proofsteps}
\item Let current accrued return be \(a\), with \(m\) future local coupons.
\item Each future coupon lies in \([\ell,u]\).
\item Therefore final accumulated return lies in
\[
a+\sum_{j=1}^m c_j \in [a+m\ell,\;a+mu].
\]
\item If \(a+mu\) is below the global floor, the payoff is exactly the global floor.
\item If \(a+m\ell\) is above the global cap, the payoff is exactly the global cap.
\item These states need no regression and no simulation.
\end{proofsteps}

\chapter{Line-by-Line Proofs: PCHIP}

\section{Unique Hermite cubic}
\begin{proofsteps}
\item A cubic has four coefficients.
\item Endpoint value constraints give two linear equations.
\item Endpoint slope constraints give two more linear equations.
\item The Hermite basis constructs one solution:
\[
H(t)=y_i h_{00}(t)+h_i d_i h_{10}(t)+y_{i+1}h_{01}(t)+h_i d_{i+1}h_{11}(t).
\]
\item If two cubics satisfied all four conditions, their difference would have double roots at both endpoints.
\item A nonzero cubic cannot have four roots counting multiplicity, so the solution is unique.
\end{proofsteps}

\section{Why PCHIP avoids artificial extrema}
\begin{proofsteps}
\item Let \(\delta_i=(y_{i+1}-y_i)/h_i\).
\item If neighboring secants change sign, PCHIP sets the node slope to zero.
\item If neighboring secants have the same sign, PCHIP uses the weighted harmonic mean
\[
d_i=\frac{w_1+w_2}{w_1/\delta_{i-1}+w_2/\delta_i}.
\]
\item This slope has the same sign as the neighboring secants and lies between their magnitudes.
\item On a monotone interval, the derivative of the Hermite cubic is a quadratic whose normalized endpoint slopes satisfy the Fritsch-Carlson monotonicity inequalities.
\item Therefore the derivative does not change sign inside the interval.
\item Hence the interpolant cannot create a new local maximum or minimum between two monotone data points.
\end{proofsteps}

\chapter{Line-by-Line Proofs: Sparse Chebyshev Ridge}

\section{Chebyshev boundedness}
\begin{proofsteps}
\item Every \(x\in[-1,1]\) can be written as \(x=\cos\theta\).
\item Define \(T_n(x)=\cos(n\theta)\).
\item Since \(|\cos(n\theta)|\le1\),
\[
|T_n(x)|\le1,\qquad x\in[-1,1].
\]
\item Scaling state variables to \([-1,1]\) keeps basis columns numerically bounded.
\end{proofsteps}

\section{Ridge normal equation}
\begin{proofsteps}
\item Objective:
\[
J(\beta)=\|A\beta-y\|_2^2+\lambda\|\Gamma\beta\|_2^2.
\]
\item Differentiate:
\[
\nabla J(\beta)=2A^\top(A\beta-y)+2\lambda\Gamma^\top\Gamma\beta.
\]
\item Set the gradient to zero:
\[
A^\top A\beta-A^\top y+\lambda\Gamma^\top\Gamma\beta=0.
\]
\item Rearrange:
\[
(A^\top A+\lambda\Gamma^\top\Gamma)\beta=A^\top y.
\]
\item The left matrix is positive semidefinite, and ridge makes unsupported directions stable.
\end{proofsteps}

\chapter{Line-by-Line Proofs: Bounded Targets and Metrics}

\section{Logit bounded target}
\begin{proofsteps}
\item Suppose the price must lie in \([L,U]\).
\item Fit an unconstrained number \(y\) and map it through the logistic function:
\[
u(y)=\frac{1}{1+e^{-y}}.
\]
\item Because \(e^{-y}>0\), \(0<u(y)<1\).
\item Rescale:
\[
\widehat V=L+(U-L)u(y).
\]
\item Therefore \(L<\widehat V<U\).
\end{proofsteps}

\section{Relative-error floor}
\begin{proofsteps}
\item Define
\[
e_{\rm rel}=\frac{|\widehat V-V|}{\max(|V|,\eta)},\qquad \eta>0.
\]
\item If \(|V|\ge\eta\), this is ordinary relative error.
\item If \(|V|<\eta\), then
\[
e_{\rm rel}=\frac{|\widehat V-V|}{\eta}.
\]
\item Thus a nearly zero benchmark cannot create an infinite percentage error.
\end{proofsteps}

\backmatter
\chapter{References}
\begin{itemize}
  \item Black, F. and Scholes, M. (1973). The Pricing of Options and Corporate Liabilities. \emph{Journal of Political Economy}.
  \item Fritsch, F. N. and Carlson, R. E. (1980). Monotone Piecewise Cubic Interpolation. \emph{SIAM Journal on Numerical Analysis}.
  \item Fritsch, F. N. and Butland, J. (1984). A Method for Constructing Local Monotone Piecewise Cubic Interpolants. \emph{SIAM Journal on Scientific and Statistical Computing}.
  \item Hoerl, A. E. and Kennard, R. W. (1970). Ridge Regression: Biased Estimation for Nonorthogonal Problems. \emph{Technometrics}.
  \item Glasserman, P. (2004). \emph{Monte Carlo Methods in Financial Engineering}. Springer.
  \item Glasserman, P. and Staum, J. (2001). Conditioning on One-Step Survival for Barrier Option Simulations. \emph{Operations Research}.
  \item Broadie, M., Glasserman, P., and Kou, S. G. (1997). A Continuity Correction for Discrete Barrier Options. \emph{Mathematical Finance}.
\end{itemize}

\end{document}
